\chapter{Strongyloides Antibody Test}

\section*{What Is This Test?}
Strongyloides antibody testing detects antibodies to \textit{Strongyloides stercoralis}, an intestinal nematode capable of lifelong autoinfection. Untreated infection can progress to fatal hyperinfection or disseminated disease after corticosteroids, transplantation, or other immunosuppression.

\section*{Alternative Names}
Strongyloides IgG; Strongyloidiasis Serology; Strongyloides Antibody; Strongyloides ELISA

\section*{How It Is Performed}
Serum is tested, commonly by an enzyme immunoassay for Strongyloides-specific IgG. Serology is generally more sensitive for screening than a single routine stool examination. Stool concentration, culture, or PCR may be used for confirmation or follow-up in selected cases.

\section*{Why It Is Ordered}
\begin{itemize}
  \item Screening people from or with prolonged travel to endemic areas before corticosteroids or other immunosuppression
  \item Evaluation of unexplained eosinophilia, recurrent gastrointestinal symptoms, or compatible pulmonary disease
  \item Investigation of suspected hyperinfection or disseminated strongyloidiasis with urgent specialist testing
\end{itemize}

\section*{Interpretation and Limitations}
A positive result supports infection but can cross-react with some other helminths and may remain positive after treatment. A negative result does not completely exclude infection, particularly in immunocompromised patients. Stool testing is insensitive when performed only once; clinicians may use serial or specialized stool testing and should not delay treatment when hyperinfection is strongly suspected.
\section*{Regulatory Status and Authorization Date}
Regulatory status applies to the specific assay, instrument, manufacturer, and intended use rather than to the test name alone. No single approval date can be assigned to this test category. For a commercial assay, verify the current product in the relevant regulator's database: the U.S. Food and Drug Administration (FDA) clearance, approval, or authorization; India's Central Drugs Standard Control Organisation (CDSCO) licence or permission; the European Union In Vitro Diagnostic Regulation (EU IVDR) CE certificate issued through the applicable conformity-assessment route; or the United Kingdom Medicines and Healthcare products Regulatory Agency (MHRA) registration and UKCA/CE status. Laboratory-developed tests may not have an individual FDA, CDSCO, EU IVDR, or MHRA approval date.
\clearpage
